\section{The ZF(C) schemas, and PA induction as the contrast}\label{sec:zf}
% Section 5: the ZF(C) schemas (Question 2) and PA induction as the contrast.
% Sources: research/tracks/cases/notes-final.md (Part 2), research/prior/induction/*-notes.md and referee reports,
% research/tracks/single/notes-final.md (Cor D.3, Thm B), research/tracks/experiments/notes-final.md (E2).
% Review revision (research/paper-review): prop:zf:false (ReplS parenthetical restricted to anchor data),
% opening overclaim fixed, (R) added to prop:zf:russell, closure caveat in Conventions; length cuts: def:zf:enc deleted,
% thm:zf:multi with its events and examples, prop:zf:indgen, prop:zf:sub and several proof ideas moved to app-zfc.tex
% (labels kept).
% Local macros (prefix zf); \providecommand so that a later move to preamble.tex does not clash.
\providecommand{\zfN}[1]{\textup{(N$_{#1}$)}}
\providecommand{\zfR}[1]{\textup{(R$_{#1}$)}}
\providecommand{\zfB}{\textup{(B)}}
\providecommand{\zfDpat}[1]{\textup{(D$^{\mathrm{PAT}}_{#1}$)}}
\providecommand{\zfDdt}[1]{\textup{(D$^{\mathrm{DT}}_{#1}$)}}
\providecommand{\zfDroot}[1]{\textup{(D$^{\mathrm{root}}_{#1}$)}}
\providecommand{\zfLinf}{L_{\infty}}

Question~2 asks whether the method developed for PA induction also learns each axiom schema of ZFC from a
collection of its instances. The answer is yes, and more easily than for induction. Every ZFC axiom other than
Separation and Replacement is a single sentence, its own anchor (\Cref{prop:zf:ground}). Every ZF schema, in each
textbook formulation considered here, is a higher-order pattern \citep{miller1991logic}: its one metavariable is
applied only to distinct bound variables. Two instances are an anchor iff their bodies have different main symbols and
jointly use every argument place, in every class between $\PAT$ and $\SO$ (\Cref{thm:zf:anchor}), and pattern
anti-unification recovers each schema from any anchor (\Cref{cor:zf:patlgg}). Plain first-order anti-unification works
iff the metavariable is applied to literally the same argument tuple at every occurrence, which depends on the
formulation and on how bound variables are written (\Cref{thm:zf:classify}). Where it fails with textbook names or de
Bruijn indices, no finite union of guarded first-order schemas covers the schema without a false member
(\Cref{thm:zf:nounion}), with one exception: spelled-out Replacement in the named encoding with textbook names, whose
first-order lgg of anchor data is a truth-sound over-generalization (\Cref{prop:zf:coll}). PA induction is the contrast. It is not a pattern, and first-order and pattern anti-unification
are both unsound on it. The determinate class $\DT$ learns it, and there the anchor condition is again root variation
plus non-vacuity (\Cref{thm:zf:indanchor}); without determinacy a stronger condition is needed
(\Cref{prop:zf:indso}). The difference is structural: $L_\in$ has no function symbols, so an instance of a ZF schema
only renames variables inside the body, while induction substitutes the terms $0$ and $Sx$.

\Cref{sec:zf:forms} fixes formulations and encodings, \Cref{sec:zf:fo,sec:zf:fofail} treat first-order learners,
\Cref{sec:zf:pattern} pattern targets, \Cref{sec:zf:pa} induction, and \Cref{sec:zf:summary} collects the answer.
Proofs, several metavariables and the computational evidence are in \Cref{app:zf}; untagged mixtures in \Cref{sec:many}.

\paragraph{Conventions.} Data are instances in closure-normal form (\Cref{sec:setting:syntax}): the textbook outer
prefix $\forall\bar w$ becomes parameters, while the frame's own variables stay bound. Anchors, closures and the
cautious verifier are as in \Cref{def:setting:protocol}; first-order lggs, guards and the most-specific-guard rule as
in \Cref{sec:setting:fo}. The ZF targets lie in $\PAT$ and $T_{\Ind}$ in $\DT$. Over a class $H$ that contains the
target, the cautious verifier is sound on all data and exact once they contain an anchor; over a class that does not,
an anchor can make it accept the larger closure $\mathrm{cl}_H(\inst(T^*))$ (\Cref{lem:setting:closed},
\Cref{prop:zf:indgen}(b)). Every first-order schema covering a finite set $P$ is at least
as general as $\lgg(P)$ \citep{plotkin1970note,reynolds1970transformational}, which is why first-order failures are
stated for lggs. Let $\top:=\forall w\,(w=w)$ and $\bot:=\neg\top$.

% ---------------------------------------------------------------------------------------------------------------
\subsection{Formulations and encodings}\label{sec:zf:forms}

Besides Separation and Replacement we treat Collection and $\in$-induction, which are theorem schemas of ZF
\status{known}: Collection via ranks or reflection (see e.g.\ \citealp{levy1979basic}; we did not check the exact
location), $\in$-induction via transitive closures and Foundation. Conversely, $\in$-induction with $\varphi(x):=x\notin
S$ proves Foundation for $S$. We include both because texts and libraries state them as schemas.
\Cref{tab:zf:forms} lists the templates. The Kunen- and Jech-style wordings follow \citet[ch.~I]{kunen1980set} and
\citet[ch.~1]{jech2003set} from memory; an independent referee's recollection agrees, but we did not re-check them
against the books. As far as we know Kunen treats $\exists!$ as an abbreviation, so $\ReplU$, with $\exists!$ as a
primitive binder, is our encoding choice.

\begin{table}[t]
\centering\small
\renewcommand{\arraystretch}{1.12}
\begin{tabularx}{\textwidth}{@{}l>{\raggedright\arraybackslash}X l@{}}
\toprule
name & template $T^*$ (closure-normal form) & arguments of $P$\\
\midrule
$\Sep$ (Kunen) & $\forall z\,\exists y\,\forall x\,(x\in y\leftrightarrow x\in z\wedge P(x,z))$ & $(x,z)$\\
$\Sep_{\mathrm J}$ (Jech) & $\forall X\,\exists Y\,\forall u\,(u\in Y\leftrightarrow u\in X\wedge P(u))$ & $(u)$\\
$\EInd$ & $\forall x\,(\forall y\,(y\in x\to P(y))\to P(x))\to\forall x\,P(x)$ & $(y)$, $(x)$, $(x)$\\
$\Coll$ & $\forall A\,(\forall x\,(x\in A\to\exists y\,P(x,y,A))\to\exists Y\,\forall x\,(x\in A\to\exists y\,(y\in
  Y\wedge P(x,y,A))))$ & $(x,y,A)$ twice\\
$\ReplU$ & as $\Coll$, with $\exists!y$ in place of the first $\exists y$ & $(x,y,A)$ twice\\
$\ReplS$ & $\forall A\,(\forall x\,(x\in A\to\exists y\,(P(x,y,A)\wedge\forall u\,(P(x,u,A)\to u=y)))\to\exists
  Y\,\forall x\,(x\in A\to\exists y\,(y\in Y\wedge P(x,y,A))))$ & $(x,y,A)$, $(x,u,A)$, $(x,y,A)$\\
$\ReplJ$ (Jech) & $\forall x\,\forall y\,\forall u\,(P(x,y)\wedge P(x,u)\to y=u)\to\forall X\,\exists Y\,\forall
  y\,(y\in Y\leftrightarrow\exists x\,(x\in X\wedge P(x,y)))$ & $(x,y)$, $(x,u)$, $(x,y)$\\
$\Found$ & $\forall S\,(\exists x\,x\in S\to\exists x\,(x\in S\wedge\forall y\,(y\in x\to\neg y\in S)))$ & ground\\
\bottomrule
\end{tabularx}
\caption[The ZF schemas as templates]{The ZF schemas as templates with one formula metavariable $P$ (closure-normal
form). Plugging a body $\lambda\bar z.\varphi$ whose free variables are holes and parameters gives a sentence.}
\label{tab:zf:forms}
\end{table}

\begin{proposition}[Single axioms]\label{prop:zf:ground}
\status{proved}\src{cases \S2.1; single Cor D.3}
Each ZFC axiom other than the two schemas, Choice included, is a single sentence $s$. Its target $\{s\}$ lies in
every class, $\{s\}$ is an anchor in every class, and matching is equality. \emph{Proof:} every template covering
$\{s\}$ contains the whole target.
\end{proposition}

Bound variables are encoded by names (Nm) or by de Bruijn indices treated as constants (dB) (\Cref{sec:setting:fo}),
or, by default, under the $\lambda$ convention (\Cref{sec:setting:templates}), where capture cannot occur and freshness
is built in. Named data use the textbook names of \Cref{tab:zf:forms} unless said otherwise; under
\emph{$\alpha$-variation} the frame's binders get arbitrary distinct names.

% ---------------------------------------------------------------------------------------------------------------
\subsection{Which ZF schemas are first-order patterns}\label{sec:zf:fo}

Let $T^*=F[P(\bar a_1),\dots,P(\bar a_m)]$ be a pattern template over $L_\in$ with a metavariable-free,
parameter-free frame $F$ and one metavariable $P$ of arity $n$, applied at its $k$-th occurrence to distinct bound
variables $\bar a_k$. Let $N_k$ be the tuple of \emph{names} of $\bar a_k$ (Nm) and $I_k$ the tuple of
their \emph{indices} at occurrence $k$ (dB).

\begin{theorem}[Classification]\label{thm:zf:classify}
\status{proved}\src{cases Thm B1}
(a) $T^*\in\PAT\subseteq\DT$.
(b) With textbook names, $\inst(T^*)$ is the instance set of a single first-order schema with freshness guards iff
$N_1=\dots=N_m$. The schema is then $F[Z,\dots,Z]$, with $\mathrm{fresh}(v,Z)$ for each frame-bound name $v$ not in
$N_1$.
(c) In de Bruijn notation the same holds with $I_1=\dots=I_m$ and the guard $\mathrm{loose}(Z)\subseteq I_1$.
Well-formedness implies this guard iff $I_1=\{0,\dots,d-1\}$, where $d$ is the least binder depth of an occurrence.
\end{theorem}
\begin{proof}[Proof idea]
If the tuples are equal, every occurrence reads the same text. If two differ, the lgg of the instances with bodies
$\bigwedge_i(z_i=z_i)$ and its negation decouples them; filling one with the closed $\top$ and the others with $\bot$
respects every freshness guard, but no body produces the result (\Cref{app:zf:enc}).
\end{proof}
\Cref{tab:zf:classify} (in \Cref{app:zf:enc}) applies the criterion to every formulation and encoding. It corrects
two natural expectations. $\in$-induction is a de Bruijn first-order pattern, since all its occurrences read $\#0$. And freshness is automatic only in the $\lambda$ encoding,
not in the de Bruijn first-order one.

\begin{proposition}[$\alpha$-variation]\label{prop:zf:alpha}
\status{proved}\src{cases Prop B1$\alpha$}
Under $\alpha$-variation (with name metavariables and guards $\mathrm{fresh}$ and $\mathrm{distinct}$), the
$\alpha$-variants of instances of $T^*$ form the instance set of a single guarded first-order schema iff all
occurrences of $P$ are applied to the same tuple of \emph{binders}. This is true for $\Sep$ and $\Sep_{\mathrm J}$,
and false for $\EInd$, $\Coll$, $\ReplU$, $\ReplS$ and $\ReplJ$.
\end{proposition}

\begin{proposition}[What the first-order lgg does]\label{prop:zf:folgg}
\status{proved}\src{cases Prop B2}
Suppose the bodies' main symbols are not all equal. Then the first-order lgg (either encoding, textbook names) is $F$
with a metavariable at each $P$-position, and positions $k,l$ share it iff no body uses an argument place at which
$\bar a_k$ and $\bar a_l$ differ. If all positions share one metavariable, the lgg with its most specific guards accepts
only instances of $T^*$.
\end{proposition}
So on anchor data the first-order lgg decouples exactly the occurrences that \Cref{thm:zf:classify} cannot tie.

% ---------------------------------------------------------------------------------------------------------------
\subsection{Where first-order anti-unification fails}\label{sec:zf:fofail}

\begin{proposition}[Two instances give a false lgg instance]\label{prop:zf:false}
\status{proved}\src{cases \S2.4}
In each of the following cases two genuine instances give an lgg with a false sentence instance that respects the
most specific guards and is refuted in ZF; all witnesses are in \Cref{app:zf:fofail}.
(i) $\EInd$, named. From $\EInd(x=x)$ and $\EInd(\neg x\in a)$ the lgg is $\forall x(\forall y(y\in x\to Z_1)\to
Z_2)\to\forall x\,Z_2$. Its instance $\forall x(\forall y(y\in x\to\neg y=y)\to\forall w\neg w\in x)\to\forall x\forall
w\,\neg w\in x$ has a valid antecedent and says that every set is empty.
(ii) $\ReplJ$, named or de Bruijn: the instance asserts a universal set.
(iii) $\Coll$, $\ReplU$, $\ReplS$, de Bruijn or $\alpha$-varied. The bodies $y=x$ and $\neg x=A$ decouple the
consequent's occurrence; $y=x$ at the antecedent's occurrence and $\bot$ at the others give an instance equivalent to
$\forall A\,\neg\exists x\,x\in A$.
(iv) Guard violations. Without its guard, Separation's lgg on data with \evR{} has Russell's instance
(\Cref{prop:zf:russell}). Without $\mathrm{fresh}(Y,\cdot)$, the named lggs of $\Coll$, $\ReplU$ and $\ReplS$ have the
instance with body $y=Y$, which forces $Y\in Y$.
For $\ReplS$ in the named encoding with textbook names, anchor data give no such instance: the lgg is truth-sound
(\Cref{prop:zf:coll}).
\end{proposition}

\begin{theorem}[No finite union of first-order schemas]\label{thm:zf:nounion}
\status{proved}\src{cases Lemma C1, Thm C2}
In each case below, every finite set of guarded first-order schemas whose instances include all instances of the
schema has a false instance that is refutable in ZF: (i) $\EInd$, named; (ii) $\ReplJ$, both encodings; (iii)
$\Coll$, $\ReplU$, $\ReplS$, de Bruijn; (iv) $\Coll$, $\ReplU$, $\ReplS$, named under $\alpha$-variation.
\end{theorem}
\begin{proof}[Proof idea]
Suppose a schema covers a sentence and has fewer nodes than the depth of one of its leaves $\ell$. Then it has a
formula metavariable $Z$ at a proper prefix $q$ of $\ell$. If the subterm at $q$ is unique, the schema also covers the
sentence with $\top$ or $\bot$ at $q$, and closed values respect every guard (\Cref{lem:zf:deepleaf}). Take instances
with deep bodies, such as $\EInd(\neg^{2n}\forall w\,\neg w\in x)$, and a leaf in the occurrence that cannot be tied.
For each prefix, one of $\top,\bot$ gives a false sentence (\Cref{app:zf:fofail}).
\end{proof}
With textbook names, the one non-first-order case not covered is $\ReplS$ in the named encoding. It is a genuine
exception:

\begin{proposition}[A truth-sound over-generalization]\label{prop:zf:coll}
\status{proved}\src{cases Props C3, C3$'$}
Let
\[L_S:=\forall A\,(\forall x\,(x\in A\to\exists y\,(Z_1\wedge\forall u\,(Z_2\to u=y)))\to\exists Y\,\forall x\,(x\in
A\to\exists y\,(y\in Y\wedge Z_1))).\]
This is the named lgg of $\ReplS$ data on which the bodies' roots vary and some body uses $y$.
Take it with the guards learned from anchor data: $\mathrm{fresh}(Y,Z_1)$, $\mathrm{fresh}(u,Z_1)$,
$\mathrm{fresh}(Y,Z_2)$, $\mathrm{fresh}(y,Z_2)$.
(a) Every guarded instance of $L_S$ is a theorem of ZF, and $\inst(L_S)\supsetneq\inst(\ReplS)$.
(b) For any theory $B$, all guarded instances of $L_S$ are theorems of $B$ iff $B$ proves the Collection schema.
\end{proposition}
\begin{proof}[Proof idea]
The antecedent implies $\forall x{\in}A\,\exists y\,Z_1$, and Collection gives the consequent. Conversely, $Z_1:=\varphi$
and $Z_2:=\bot$ give a formula equivalent to $\Coll(\varphi)$, since $\forall u(\bot\to u=y)$ is valid. This formula is
not an instance of $\ReplS$ (\Cref{app:zf:fofail}).
\end{proof}
So here the first-order learner is truth-sound but not exact, and no refutation can ever detect its
over-generalization. Non-anchor data can still give false instances when freshness guards apply only to formula
metavariables (\Cref{app:zf:fofail}). With Jech's image consequent the decoupling gives false instances instead, and
under $\alpha$-variation the phenomenon disappears. By (b), the lgg is truth-sound over ZF without Foundation iff that
theory proves Collection; we do not know whether it does.

\begin{proposition}[Separation: the guard and Russell's instance]\label{prop:zf:russell}
\status{proved}\src{cases Prop D1}
Let the bodies' main symbols not all be equal \evR.
(a) The unguarded lgg of Separation data (either encoding) has Russell's instance $\forall z\exists y\forall x(x\in
y\leftrightarrow x\in z\wedge\neg x\in y)$. It is refuted in pure logic from $\exists z\exists x\,x\in z$.
(b) The most specific guard rule always learns $\mathrm{fresh}(y,Z)$, since no genuine body mentions $y$. The learned
guards equal the target's iff some body uses $x$ and some body uses $z$. For $\Coll$ and $\ReplU$ (named), the learned
guards equal the target's, $\{\mathrm{fresh}(Y,Z)\}$, iff every argument place is used.
\end{proposition}
\begin{proof}
By \evR{} the lgg has a formula metavariable $Z$ at the slot. (a) $Z:=\neg x\in y$; for $x_0\in z$ and the $y$ given by
the sentence, $x_0\in y\leftrightarrow\neg x_0\in y$. (b) A freshness guard for $v$ is learned iff no datum's value
mentions $v$.
\end{proof}

% ---------------------------------------------------------------------------------------------------------------
\subsection{Pattern targets: anchors, pattern anti-unification and rates}\label{sec:zf:pattern}

Let $T^*$ be as in \Cref{sec:zf:fo} and $D=\{T^*[P:=\lambda\bar z.\varphi_j]:j\le N\}$. Let \evR{} say that the main
symbols of the bodies $\varphi_j$ are not all equal, and \zfN{i} ($i\le n$) that some $\varphi_j$ has the hole $z_i$
free.

\begin{theorem}[Anchors for the ZF schemas]\label{thm:zf:anchor}
\status{proved}\src{cases Thm E; single Cor D.3}
For every class $H$ with $\PAT\subseteq H\subseteq\SO$, a finite $D\subseteq\inst(T^*)$ is an anchor in $H$ iff
\evR{} and \zfN{i} hold for all $i\le n$. In particular one instance never suffices, and two can.
\end{theorem}
\begin{proof}[Proof idea]
\evR{} pins the rigid part of every covering template to a prefix of the frame, and \zfN{i} lets every variable bound
above an occurrence be recovered as a hole of the body; so a covering template, with a suitable body for each
metavariable, gives back $T^*$. ``Only if'': a root-specialized or argument-dropping pattern covers non-anchor data
(\Cref{app:zf:anchor}).
\end{proof}
Minimal anchors, as bodies, are: $x\in z$ and $\neg(x=a)$ for $\Sep$; $x\in x$ and $\neg(x=x)$ for $\EInd$; $y=x$ and
$\neg(x=A)$ for $\Coll$, $\ReplU$, $\ReplS$; $x\in y$ and $\neg(x=a)$ for $\ReplJ$. If a community never lets $\varphi$
mention the bounding set $z$, the learner stays at the sound specialization $P(x)$, which is Jech's form. Since
$L_\in$ has no function symbols, for one metavariable the ``shielding'' that separates $\SO$ from $\DT$ for induction
(\Cref{prop:zf:indso}) and for universal axioms over arithmetic (\Cref{prop:univ:dt}(c)) cannot occur. For $H=\DT$ the
theorem is also a case of \Cref{cor:single:special}(c).

\paragraph{Several metavariables.} ZF needs one metavariable. With several, as in untagged mixtures
(\Cref{sec:many}), a Plotkin-(D)-type event appears that differs between $\PAT$, $\DT$ and $\SO$, so that $\SO$ and
$\DT$ anchors differ even over $L_\in$ (\Cref{thm:zf:multi}); the exact $\SO$ condition is open (\Cref{conj:zf:root}).

\begin{corollary}[Pattern anti-unification recovers every ZF schema]\label{cor:zf:patlgg}
\status{proved}\src{cases Cor F, F$'$}
The least general pattern generalization $L_{\mathrm{pat}}(D)$ exists and is unique
\citep{pfenning1991unification,baumgartner2017higher}. For $D\subseteq\inst(T^*)$, $L_{\mathrm{pat}}(D)\preceq T^*$,
and $L_{\mathrm{pat}}(D)\equiv T^*$ iff \evR{} and every \zfN{i} hold. With several metavariables,
$L_{\mathrm{pat}}(D)\equiv T^*$ iff \zfR{P}, every \zfN{P,i} and \zfDpat{PQ} for all $P\ne Q$ hold
(\Cref{thm:zf:multi}).
\end{corollary}
\begin{proof}
$T^*\in\PAT$ covers $D$. Under the events, \Cref{thm:zf:anchor} (resp.\ \ref{thm:zf:multi}) for $H=\PAT$ gives
$L_{\mathrm{pat}}(D)\gen T^*$. Otherwise its $\PAT$ witnesses lie above $L_{\mathrm{pat}}(D)$ and miss an instance.
\end{proof}
The merge rule must identify columns that are equal up to a permutation of the arguments: Replacement's $P(x,y)$ and
$P(x,u)$ columns agree only up to $y\leftrightarrow u$.

\begin{theorem}[Rates]\label{thm:zf:rates}
\status{proved}\src{cases Thm G}
Let the bodies be i.i.d.\ from $\mu$. Let $C_{F,I}$ be the set of bodies with root $F$ (no condition if $F={*}$) in
which no $z_i$, $i\in I$, is free. Then
\[\Pr[D_N\text{ is not an anchor}]=\sum_{(F,I)\ne({*},\emptyset)}(-1)^{[F\ne{*}]+|I|+1}\mu(C_{F,I})^N\;\le\;\sum_f
p_f^N+\sum_i(1-q_i)^N,\]
where $p_f=\mu(\mathrm{root}=f)$ and $q_i=\mu(z_i\text{ free})$.
\end{theorem}
\begin{proof}
Inclusion--exclusion over the events ``all roots are $f$'' (pairwise disjoint) and ``$z_i$ is never free''.
\end{proof}

% ---------------------------------------------------------------------------------------------------------------
\subsection{PA induction: the contrast}\label{sec:zf:pa}

A \emph{motive} is a formula $\varphi$ of $L_A$ with $\FV(\varphi)\subseteq\{x\}$, and
$\Ind(\varphi)=\varphi(0)\wedge\forall x(\varphi\to\varphi[Sx/x])\to\forall x\varphi$. The target is
$\inst(T_{\Ind})$ with $T_{\Ind}=P(0)\wedge\forall x(P(x)\to P(Sx))\to\forall x\,P(x)$. Since $P(0)$ and $P(Sx)$ are
not pattern occurrences, $T_{\Ind}\in\DT\setminus\PAT$. Here \evR{} says that the motives' main symbols are not all
equal, and \evN{} that some motive has $x$ free. The results are from the earlier analysis, for parameter-free motives
unless said otherwise; proofs are in \Cref{app:zf:pa}.

\paragraph{First-order learners.} A first-order schema is \emph{truth-sound} if its sentence instances are true in
$\N$. The false instances below are closed and capture-free, so the results hold in the named and in the de Bruijn
encoding, with fixed or varying bound-variable names.

\begin{proposition}[Two raw instances suffice for unsoundness]\label{prop:zf:indtwo}
\status{proved}\src{prior induction B1, B3}
(a) $\lgg(\Ind(x+0=x),\Ind(0+x=x))=(0+0=0)\wedge\forall x(z_1+z_2=x\to z_3+z_4=Sx)\to\forall x(z_1+z_2=x)$. It has
the false instance $(0+0=0)\wedge\forall x(0+0=x\to0+S0=Sx)\to\forall x(0+0=x)$. One raw instance never forces
unsoundness, and some pairs, such as $\Ind(x=x)$ and $\Ind(Sx=Sx)$, have a sound lgg.
(b) If \evR{} and \evN{} hold, then $\lgg(D)=\zfLinf:=Z_1\wedge\forall x(Z_2\to Z_3)\to\forall x\,Z_2$, which has the
false instance $(0=0)\wedge\forall x(0=S0\to0=S0)\to\forall x(0=S0)$. For i.i.d.\ motives,
$\Pr[\lgg(D_N)\ne\zfLinf]\le p_{\max}^{N-1}+(1-q)^N$.
\end{proposition}
So the diversity that identifies $T_{\Ind}$ among determinate templates (\Cref{thm:zf:indanchor}) destroys the
first-order lgg, at the same coupon-collector rate.

\begin{theorem}[No finite union of first-order schemas]\label{thm:zf:indunion}
\status{proved}\src{prior induction B2}
Let $g_0(t)=t$, $g_{n+1}(t)=0+g_n(t)$ and $F'=\{\Ind(g_n(x)=x):n\ge0\}$. Every first-order schema covering two members
of $F'$ has a false instance that is refutable in $\Q$. Hence $k$ truth-sound first-order schemas cover at most $k$
members of $F'$, and no finite union of first-order schemas covers the raw induction instances (or those of open
induction) truth-soundly.
\end{theorem}
The occurrence and freshness guards learnable from $F'$ do not help (\Cref{app:zf:pa}); a guard ``is an induction
instance'' would, which is what the Sub encoding below implements. Worse, some unsound lggs are never caught by
evaluating closed sentences.

\begin{theorem}[An unsound lgg that closed truths never refute]\label{thm:zf:K0}
\status{proved}\src{prior induction referee R1}
Let $K_0:=(z_1+0=z_1)\wedge\forall x(z_1+x=z_2\to z_1+Sx=Sz_2)\to\forall x(z_1+x=z_2)$, the lgg of $\Ind(0+x=x)$ and
$\Ind(S0+x=Sx)$. Its instance $(0+0=0)\wedge\forall x(0+x=0\to0+Sx=S0)\to\forall x(0+x=0)$ is false. Yet the true
closed quantifier-free sentences together with all sentence instances of $K_0$ are consistent. So no sound derivation
from true closed quantifier-free premises (or $\Delta_0$ ones, with $<$ primitive) refutes an instance of $K_0$.
\end{theorem}
$K_0$ is refuted once the recursion axiom Q5, $\forall x\forall y\,(x+Sy=S(x+y))$, is a designated premise \src{prior
induction B5(d)}.

\paragraph{Pattern anti-unification and determinate templates.}

\begin{proposition}[Pattern anti-unification is unsound on induction]\label{prop:zf:indpat}
\status{proved}\src{prior induction C9}
Under \evR{} and \evN, the least general pattern generalization of $D\subseteq\inst(T_{\Ind})$ is
$T_{\mathrm{pat}}:=Z\wedge\forall x(P(x)\to Q(x))\to\forall x\,P(x)$. It has the false non-instance
$(0=0)\wedge\forall x(x=0\to x=0)\to\forall x(x=0)$.
\end{proposition}

\begin{table}[t]
\centering\small
\setlength{\tabcolsep}{4pt}
\begin{tabularx}{\textwidth}{@{}l c c c c c >{\raggedright\arraybackslash}X@{}}
\toprule
& \multicolumn{3}{c}{first-order} & & & \\
\cmidrule(lr){2-4}
target & named & dB & named, $\alpha$ & $\PAT$ & $\DT$ & anchor condition\\
\midrule
ground axioms (incl.\ $\Found$, $\AC$) & yes & yes & yes & yes & yes & the sentence itself\\
$\Sep$ & g & g & g & yes & yes & \evR, \zfN{x}, \zfN{z}\\
$\Sep_{\mathrm J}$ & g & g & g & yes & yes & \evR, \zfN{u}\\
$\EInd$ & no$^{\dagger}$ & yes & no & yes & yes & \evR, \zfN{x}\\
$\Coll$, $\ReplU$ & g & no$^{\dagger}$ & no$^{\dagger}$ & yes & yes & \evR, \zfN{x}, \zfN{y}, \zfN{A}\\
$\ReplS$ & no$^{\ddagger}$ & no$^{\dagger}$ & no$^{\dagger}$ & yes & yes & \evR, \zfN{x}, \zfN{y}, \zfN{A}\\
$\ReplJ$ & no$^{\dagger}$ & no$^{\dagger}$ & no & yes & yes & \evR, \zfN{x}, \zfN{y}\\
\midrule
$\Ind$, raw & no$^{\dagger}$ & no$^{\dagger}$ & no$^{\dagger}$ & no$^{\S}$ & yes & $\DT$: \evR, \evN; $\SO$: \evR,
  \zfB\\
$\Ind$, Sub-encoded & yes & -- & -- & -- & -- & first-order lgg: \evR, \evN\\
$\IndEq$ & g & -- & -- & -- & -- & first-order lgg: \evR\\
\bottomrule
\end{tabularx}
\caption[Summary of the section]{All targets of this section. ``yes'': a first-order pattern (resp.\ in the class).
``g'': a first-order pattern with guards (e.g.\ $\mathrm{fresh}(y,Z)$ for $\Sep$, $\mathrm{fresh}(Y,Z)$ for $\Coll$,
$y\notin\FV(P)$ for $\IndEq$; all: \Cref{tab:zf:classify}), with false instances without them.
``no$^\dagger$'': no finite union of guarded first-order schemas covers the target without a false member
(\Cref{thm:zf:nounion,thm:zf:indunion}). ``no$^\ddagger$'': not first-order, but the first-order lgg of anchor data is
truth-sound iff the base theory proves Collection (\Cref{prop:zf:coll}). ``no$^\S$'': not a pattern, and the pattern
lgg is unsound (\Cref{prop:zf:indpat}). ``--'': not analysed. For the ZF schemas the anchor condition is the same in
every class between $\PAT$ and $\SO$ (\Cref{thm:zf:anchor}). Two instances suffice in every row except the ground
axioms, where one does.}
\label{tab:zf:summary}
\end{table}

For the size-bounded classes $\DT_s$ (templates of size $\le s$), $|T|$ counts symbols, an occurrence $M(\bar t)$
counting $1+\sum_i|t_i|$; so $|T_{\Ind}|=14$.

\begin{theorem}[Anchors for induction in $\DT$]\label{thm:zf:indanchor}
\status{proved}\src{prior induction C3; single Cor D.3}
A finite $D\subseteq\inst(T_{\Ind})$ is an anchor in $\DT$ iff \evR{} and \evN. The same holds in the size-bounded
classes $\DT_s$ for $s\ge7$, and, by \Cref{cor:single:special}(c), in closure-normal form with parameters. Two
instances suffice; one never does.
\end{theorem}
Up to renaming exactly $26$ templates of $\DT$ contain $\inst(T_{\Ind})$, and on every anchor they are the covering
templates, so $\lgg_{\DT}(D)=T_{\Ind}$; $T_{\mathrm{pat}}$ is the least of them without non-pattern occurrences. In
$\DT_s$, $\inst(T_{\Ind})$ is closed iff $s\ge12$; for $7\le s\le11$ the cautious verifier accepts a false sentence on
every anchor (\Cref{prop:zf:indgen}, an instance of \Cref{lem:setting:closed}).

\begin{proposition}[Without determinacy]\label{prop:zf:indso}
\status{proved}\src{prior induction C7}
Call a free occurrence of $x$ \emph{unshielded} if it lies inside no proper term $t\ne x$ with
$\FV(t)\subseteq\{x\}$, and let \zfB{} say that some motive has one. In $\SO$, $D$ is an anchor iff \evR{} and \zfB.
For example, $\{\Ind(Sx=S0),\Ind(\neg Sx=0)\}$ satisfies \evR{} and \evN, but $Z\to\forall x\,P(Sx)$ covers it.
\end{proposition}

\paragraph{Encodings that make induction first-order.}
Recording each step with the decidable judgment $\mathrm{Sub}(\varphi,x,t,\psi)$ (``$\psi=\varphi[t/x]$'') makes
induction a first-order pattern whose lgg recovers it under the same events \evR+\evN\ (\Cref{prop:zf:sub}); so the
second-order learner on raw sentences loses nothing against a first-order learner on Sub-annotated data.

\begin{proposition}[PA is $\Q$ plus one first-order pattern]\label{prop:zf:indeq}
\status{proved}\src{prior induction B6}
Let $\IndEq(P):=\forall x(x=0\to P)\wedge\forall y(\forall x(x=y\to P)\to\forall x(x=Sy\to P))\to\forall x\,P$, with
$P$ a 0-ary metavariable and $x,y$ fixed \citep{tarski1965simplified}.
(a) First-order logic proves $\IndEq(\varphi)\leftrightarrow\Ind(\varphi)$ for every motive $\varphi$, up to renaming a
bound variable. With parameters, the guard $y\notin\FV(P)$ is needed.
(b) $\lgg(\IndEq(\psi_1),\dots,\IndEq(\psi_N))=\IndEq(\lgg(\psi_1,\dots,\psi_N))$, which is $\IndEq(P)$ iff \evR{} holds.
(c) $\Q+\inst(\IndEq)\vdash\PA$.
\end{proposition}
So the obstruction in \Cref{thm:zf:indunion} is the meta-level substitution in $\varphi(0)$ and $\varphi(Sx)$, not
induction, and it differs from the non-finite axiomatizability of PA \citep{rylln1952axiomatizability}. Proof
assistants record or recompute the motive (Lean, Metamath: \Cref{app:zf:pa}); in $\DT$ the motive is the unique
matcher of $T_{\Ind}$.

% ---------------------------------------------------------------------------------------------------------------
\subsection{Summary}\label{sec:zf:summary}

The answer to Question~2 is yes (\Cref{tab:zf:summary}). All ZF schemas lie in the pattern fragment; pattern
anti-unification and the $\DT$ learner learn each from two suitable instances and are sound at all times. A schema is a
first-order pattern iff its metavariable's argument tuple is literally the same at every occurrence in the chosen
syntax (\Cref{thm:zf:classify,prop:zf:alpha}); otherwise, with textbook names or de Bruijn indices, every finite union
of guarded first-order schemas covering it has a false member, except for spelled-out Replacement with textbook names
(\Cref{thm:zf:nounion,prop:zf:coll}). PA induction lies
outside the pattern fragment; $\DT$ handles it with the anchor events of the first-order learner on Sub-annotated data,
and without determinacy the anchors need the stronger event \zfB. The experiments agree: on Separation,
$\in$-induction, a Replacement variant and induction the implemented $\DTF$ learner was exact from $2$--$3$ tagged
instances in most runs and accepted no non-instance in $927$ probes, while the pattern lgg failed on induction and
unguarded first-order lggs failed except on $\EInd$ in de Bruijn form (\Cref{tab:exp:e2,tab:zf:classify}). Open
problems: \Cref{sec:disc:open}.
